% generated by roam.py handbook export -- не редактировать, править заметки
% 20260918102313-внутреннее_описание_выпуклои_оболочки.org (body b16acfedea0a)

\begin{theorem}[{Внутреннее описание выпуклой оболочки}]\label[theorem]{note:3A5B0E29-AED4-409E-A823-956F3A3957F7}
  Пусть $V$~--- \hyperref[note:1B626AA4-35C8-4C7C-9D3D-D8C26D98F40C]{векторное пространство}, $M \subset V$.
  Тогда
  \[
  \mathrm{conv}(M) = \left\{ \sum\limits_{k=1}^n \lambda_k x_k\  \left| \  x_k \in M, \sum\limits_{k=1}^n \lambda_k = 1, \lambda_k \ge 0\right. \right\}
  \]

  То есть \hyperref[note:1AC657F9-2A3D-47D1-A71D-5B813D961260]{выпуклая оболочка} множества~--- это всевозможные \hyperref[note:8583F986-A6E8-4BBB-B2FB-BA0E9BBE58BF]{выпуклые комбинации} его элементов.
\end{theorem}

\begin{proof}
  Пусть
  \[
  S = \left\{ \sum\limits_{k=1}^n \lambda_k x_k\  \left| \  x_k \in M, \sum\limits_{k=1}^n \lambda_k = 1, \lambda_k \ge 0\right. \right\}
  \]

  \proofpart{Сперва докажем, что $\mathrm{conv}(M) \supset S$}
  По определению \hyperref[note:1AC657F9-2A3D-47D1-A71D-5B813D961260]{выпуклой оболочки}:
  \[
  \mathrm{conv}(M) = \bigcap_Q \left\{ Q \subset V, Q \text{ -- выпуклое и } M \subset Q \right\}
  \]

  А из \hyperref[note:42DFF6DC-CC9D-40AA-9D6F-218DCF613833]{леммы} мы знаем, что \hyperref[note:316B5BFA-8B9F-4267-9086-D6593AAB5F9B]{выпуклое множество} содержит все свои выпуклые комбинации.

  Получается, что $x_k \in M \Rightarrow x_k \in Q \supset M$, для произвольного выпуклого $Q \supset M$.
  Тогда \hyperref[note:8583F986-A6E8-4BBB-B2FB-BA0E9BBE58BF]{выпуклая комбинация} $x_k \in M$, тоже принадлежит к $Q \supset M$.

  Получили, что $S \subset Q$ для любого выпуклого $Q \supset M$, следовательно $S$ лежит и в их пересечении $\Rightarrow S \subset \mathrm{conv}(M)$

  \proofpart{Теперь докажем, что $S \supset \mathrm{conv}(M)$}
  Так как $\conv(M)$ это пересечение всех \hyperref[note:316B5BFA-8B9F-4267-9086-D6593AAB5F9B]{выпуклых надмножеств} $M$, то для любого
  \hyperref[note:316B5BFA-8B9F-4267-9086-D6593AAB5F9B]{выпуклого} $Q \supset M$ верно, что $Q \supset \conv(M)$.

  Следовательно, нам нужно доказать, что $S$~--- \hyperref[note:316B5BFA-8B9F-4267-9086-D6593AAB5F9B]{выпукло} и $M \subset S$.

  \proofstep{$M \subset S$}
  Так как $S$~--- все \hyperref[note:8583F986-A6E8-4BBB-B2FB-BA0E9BBE58BF]{выпуклые комбинации} $M$, то можем рассмотреть вырожденный случай \hyperref[note:8583F986-A6E8-4BBB-B2FB-BA0E9BBE58BF]{выпуклой комбинации}~--- суммы из одного элемента с $\lambda =1$.

  Тогда получим, что каждый элемент из $M$ лежит и в $S$.

  \proofstep{$S$~--- выпуклое}
  Проверим, по определению \hyperref[note:316B5BFA-8B9F-4267-9086-D6593AAB5F9B]{выпуклого множества}.

  Рассмотрим два элемента $u,v \in S$. Мы знаем, как они устроены (по определению они \hyperref[note:8583F986-A6E8-4BBB-B2FB-BA0E9BBE58BF]{выпуклые комбинации} элементов из $M$):
  \[
  u = \sum\limits_{i=1}^n \lambda_i x_i,\quad v = \sum\limits_{j=1}^m \eta_jy_j,\qquad
  \sum\limits_{i=1}^n \lambda_i = \sum\limits_{j=1}^m \eta_j = 1,\ \eta_j, \lambda_i \ge 0,\ x_i,y_j \in M
  \]

  Возьмем $\phi \in [0;1]$, рассмотрим $\phi u + (1-\phi) v$, подставим суммы
  \[
   \sum\limits_{i=1}^n (\phi \lambda_i) x_i + \sum\limits_{j=1}^m ((1-\phi)\eta_j)y_j
  \]

  Получили комбинацию из $n+m$ точек с коэффициентами $\phi\lambda_i$ и $(1-\phi)\eta_j$ соответственно.

  Коэффициенты неотрицательны; если их сумма 1~--- то это будет \hyperref[note:8583F986-A6E8-4BBB-B2FB-BA0E9BBE58BF]{выпуклой комбинацией}, проверим
  \[
  \sum\limits_{i=1}^n \phi \lambda_i + \sum\limits_{j=1}^m (1-\phi) \eta_j = \phi \cdot \sum\limits_{i=1}^n \lambda_i + (1-\phi) \sum\limits_{j=1}^m \eta_j = \phi \cdot 1 + (1-\phi) \cdot 1 = 1
  \]

  Получили, что $\sum\limits_{i=1}^n (\phi \lambda_i) x_i + \sum\limits_{j=1}^m ((1-\phi)\eta_j)y_j$~--- \hyperref[note:8583F986-A6E8-4BBB-B2FB-BA0E9BBE58BF]{выпуклая комбинация} элементов из $M$, следовательно она лежит в $S$.

  Следовательно, мы доказали, что $\forall u, v \in S\ \forall \phi \in [0; 1]\ \phi u + (1-\phi) v \in S$ \\
  Это в точности определение \hyperref[note:316B5BFA-8B9F-4267-9086-D6593AAB5F9B]{выпуклого множества}
\end{proof}
